% =====================================================================
% Numerical assessment of the Robin-residual transmission condition
% for the 3D-1D coupled problem.
%
% Drop into the DD-ROM note. Requires: amsmath, amssymb, booktabs.
%
% Add to the preamble if not already present:
%   \newcommand{\e}[1]{\!\times\!10^{#1}}
%   \newtheorem{remark}{Remark}          % or \theoremstyle{remark}
% =====================================================================

\section{Numerical assessment of the Robin-residual extraction}
\label{sec:robin-residual-assessment}

We report a controlled numerical study of the transmission condition of
eqs.~\eqref{eq:discrete_residual}--\eqref{eq:local_solve}. The conclusion is that the
scheme is \emph{exact} whenever no vessel crosses an artificial interface, and
that it \emph{fails} --- performing worse than no coupling at all --- as soon as
one does. We identify the missing term responsible and verify that supplying it
restores the structural property the extraction relies on.

\subsection{Test configuration}

To remove every confounding factor present in the full problem (curved physical
boundary, exterior-degree-of-freedom elimination, branching vasculature, eight
subdomains with three face-neighbours each) we use the smallest configuration in
which eqs.~\eqref{eq:discrete_residual}--\eqref{eq:local_solve} still carry content:

\begin{itemize}
  \item domain $\Omega = [0,1]^3$, discretised with $12^3$ cells
        ($2197$ degrees of freedom), every cell interior, so no
        exterior-dof elimination occurs and $A_i$ contains no
        unit-diagonal placeholder rows;
  \item a single straight vessel segment $\Lambda$ with two nodes, radius
        $R = 0.08$, running parallel to the $y$-axis, carrying a prescribed
        pressure drop;
  \item one artificial cut plane at $x = \tfrac12$, giving two subdomains
        $\Omega_0, \Omega_1$ sharing exactly one face.
\end{itemize}

The only quantity varied is the position of the vessel axis relative to the cut:

\begin{center}
\begin{tabular}{lll}
\toprule
label & axis position & averaging circle \\
\midrule
\textsf{offset} & $x = 0.25$ & entirely inside $\Omega_0$ \\
\textsf{near}   & $x = 0.46$ & spans the cut ($0.46 + R = 0.54 > 0.5$) \\
\bottomrule
\end{tabular}
\end{center}

\subsection{Diagnostics}

\paragraph{Interface support of the residual.}
Eq.~\eqref{eq:discrete_residual} applies the restriction $P_i^\Gamma$ to
\begin{equation}
  r_i \;=\; b_i - A_i u^*_i + \rho\, G_i^\Gamma u^*_i ,
  \label{eq:residual}
\end{equation}
discarding every entry off the artificial interface. This is legitimate only if
those entries vanish, which is precisely the \emph{Key point}: for a one-sided
local operator, $b_i - A_i u^*$ is supported on $\Gamma_i^{\mathrm{art}}$. We
therefore measure
\begin{equation}
  \eta_i \;=\;
  \frac{\bigl\| r_i|_{\Gamma_i^{\mathrm{art}}} \bigr\|_2}{\|r_i\|_2}
  \;\in\; [0,1],
  \label{eq:eta}
\end{equation}
so that $\eta_i = 1$ certifies the hypothesis and $\eta_i < 1$ quantifies the
fraction of the datum that $P_i^\Gamma$ silently destroys.

\paragraph{Error against the uncoupled baseline.}
Every error reported below is the relative $L^2$ distance between a subdomain
solution and the restriction of the global solution to that subdomain,
\begin{equation}
  e(u_i) \;=\;
  \frac{\bigl\| u^*|_{\Omega_i} - u_i \bigr\|_2}
       {\bigl\| u^*|_{\Omega_i} \bigr\|_2},
  \label{eq:relerr}
\end{equation}
averaged over the two subdomains. Here $u^*|_{\Omega_i}$ is the target --- the
converged global solution restricted to $\Omega_i$ --- and $u_i$ is whatever the
scheme under test produced there.

The column headed ``rel.\ error vs.\ raw'' in Table~\ref{tab:results} reports
\eqref{eq:relerr} for \emph{two} choices of $u_i$:
\begin{itemize}
  \item $e(u_i^{\mathrm{scheme}})$, from the local solve
        \eqref{eq:local_solve};
  \item $e(u_i^{\mathrm{raw}})$, the \emph{raw} solve, in which each
        subdomain is solved in isolation with the natural condition on its
        artificial cut planes and no transmission condition whatsoever.
\end{itemize}
The raw solve is the do-nothing baseline. It is the correct reference because
some error is unavoidable at any discretisation, so the question is not whether
a transmission condition leaves error but whether it improves on ignoring the
neighbouring subdomains altogether. A scheme with
$e(u_i^{\mathrm{scheme}}) > e(u_i^{\mathrm{raw}})$ is actively harmful.

\begin{remark}
The raw baseline is \emph{not} a single fixed number across the table: it
depends on how the local coupling operator $C_i$ is assembled, and the cross
term of \S\ref{sec:cross} requires a different assembly (see
\eqref{eq:cross} and the discussion following it). Consequently rows~2 and~3 of
Table~\ref{tab:results} have baselines $7.2\times10^{-2}$ and
$5.0\times10^{-1}$ respectively, and only the ratio
$e(u_i^{\mathrm{scheme}})/e(u_i^{\mathrm{raw}})$ is comparable between rows.
Absolute errors must never be compared across rows.
\end{remark}

\subsection{Results}

\begin{table}[htbp]
\centering
\begin{tabular}{llccccl}
\toprule
& & & \multicolumn{3}{c}{relative error \eqref{eq:relerr}} & \\
\cmidrule(lr){4-6}
vessel & scheme & $\eta_0 / \eta_1$
       & scheme & raw & ratio & \\
\midrule
\textsf{offset} & eqs.~\eqref{eq:discrete_residual}--\eqref{eq:local_solve}
  & $100 / 100$ & $2.9\e{-15}$ & $5.8\e{-1}$ & $5\e{-15}$ & exact \\
\textsf{near}   & eqs.~\eqref{eq:discrete_residual}--\eqref{eq:local_solve}
  & $38 / 86$   & $2.2\e{-1}$  & $7.2\e{-2}$ & $3.0$      & harmful \\
\textsf{near}   & \;+ cross term \eqref{eq:cross}
  & $100 / 100$ & $3.4\e{-1}$  & $5.0\e{-1}$ & $0.68$     & improves \\
\textsf{offset} & \;+ cross term \eqref{eq:cross}
  & $100 / 100$ & $2.9\e{-15}$ & $5.8\e{-1}$ & $5\e{-15}$ & unchanged \\
\bottomrule
\end{tabular}
\caption{Interface support $\eta_i$ (in \%) of the extracted residual, and the
resulting subdomain error \eqref{eq:relerr}. The \emph{ratio} column is
$e(u_i^{\mathrm{scheme}})/e(u_i^{\mathrm{raw}})$: values below $1$ mean the
transmission condition improves on solving the subdomains in isolation, values
above $1$ mean it is worse than no coupling at all. Rows~2 and~3 have different
raw baselines because the cross term requires the modified assembly of
\S\ref{sec:cross}; the ratio is therefore the only quantity comparable between
rows, and absolute errors must not be compared across them.}
\label{tab:results}
\end{table}

Row~1 establishes that the method and its implementation are correct: with no
vessel crossing the interface, eqs.~\eqref{eq:discrete_residual}--\eqref{eq:local_solve}
reproduce the global solution to machine precision, the interior part of $r_i$
being $O(10^{-18})$. Row~2 changes only the position of the vessel axis, by
$0.21$ in $x$: the interface support collapses to $38\%$ on one subdomain and
the scheme becomes three times \emph{worse} than applying no transmission
condition at all. Since nothing else differs between the two rows --- same
code, same mesh, same $\rho$ --- the failure is attributable to the vessel
crossing the interface and to nothing else.

\subsection{The missing term}
\label{sec:cross}

The local operator is
\begin{equation}
  A_i \;=\; \underbrace{AD^{00}_i}_{\text{3D diffusion}}
          \;+\; \underbrace{C_i^\top G\, C_i}_{\text{3D--1D coupling}} .
\end{equation}
A direct row-by-row comparison against the global operator, restricted to rows
whose entire global stencil lies inside $\Omega_i$, gives
$\max|AD^{00}_i - AD^{00}| = 0$ to machine precision: the diffusion block is
exactly one-sided and is never the cause.

The coupling block is not. The matrix $C$ realises the circular average
\begin{equation}
  \bar u(s) \;=\; \frac{1}{2\pi R}\oint_{\partial D(s)} u \, \mathrm{d}\theta ,
\end{equation}
so each row of $C$ reaches the 3D degrees of freedom on one circle of radius $R$.
When that circle intersects the cut plane the row splits between the two
subdomains,
\begin{equation}
  C_v \;=\; C_v^{(i)} + C_v^{(j)} ,
\end{equation}
and neither subdomain can evaluate the average alone. Restricting the global
equation to the rows owned by $\Omega_i$ gives
\begin{equation}
  AD^{00}_i u_i \;+\; C_i^\top G\, C_i u_i
  \;+\; \underbrace{C_i^\top G\, C_j u_j}_{\text{absent from \eqref{eq:local_solve}}}
  \;=\; b_i .
  \label{eq:truelocal}
\end{equation}
The third term couples $\Omega_i$ to its neighbour \emph{through the vessel
network} rather than through the flux across $\Gamma$. Since
eq.~\eqref{eq:local_solve} transports only the latter, this contribution is
unaccounted for; it is what remains in $b_i - A_i u^*$ away from the interface,
and it is distributed through the volume wherever vessels run. This is the
$62\%$ of $\|r_i\|$ that $P_i^\Gamma$ discards in row~2 of
Table~\ref{tab:results}.

We therefore augment eq.~\eqref{eq:local_solve} with
\begin{equation}
  \mathcal{X}_i \;=\; C_i^\top G \left( C u^* - C_i u_i \right),
  \label{eq:cross}
\end{equation}
written against the global $C$ so that no assumption is needed about how the
columns on $\Gamma$ are apportioned between the two subdomains. The quantity
crossing the interface is $C_v u^* - C_v^{(i)} u_i$, one scalar per straddling
vessel node --- the partial arc average on the far side --- so the exchange
remains an interface datum of the same character as the Robin trace, not a
volume field.

\paragraph{Consequence for the assembly, and for the baseline.}
Expression \eqref{eq:cross} is meaningful only if $C_v^{(i)}$ and $C_v^{(j)}$
carry the \emph{global} arc weights, so that $C_v^{(i)} + C_v^{(j)} = C_v$
reconstructs the full circular average. The straightforward local assembly does
not have this property: quadrature points falling outside $\Omega_i$ are
discarded and the partial arc is then renormalised by its own surviving
measure, which makes each fragment look like a complete average. Using
\eqref{eq:cross} therefore requires $C_i$ to be the global row with the
non-owned columns \emph{deleted} rather than renormalised.

This changes the local operator, and hence the raw baseline: a subdomain now
retains only the weight physically inside it --- between $49\%$ and $89\%$ for
a straddling vessel in this configuration --- instead of a renormalised
substitute. That is why the raw error in row~3 of Table~\ref{tab:results}
($5.0\times10^{-1}$) exceeds that of row~2 ($7.2\times10^{-2}$): the two rows
use different assemblies, and the comparison of interest is between each scheme
and the baseline computed with its own assembly.

Rows~3 and~4 of Table~\ref{tab:results} record the effect. The interface support
returns to $100\%$ on both subdomains, with interior residual $O(10^{-17})$: the
\emph{Key point} is restored exactly for a straddling vessel, and the scheme
passes from three times worse than the uncoupled baseline to $0.68$ times it.
Row~4 is a regression check --- with no vessel crossing the cut,
$\mathcal{X}_i$ vanishes identically ($\|\mathcal{X}_i\| = 0$ in floating
point) and the machine-precision result of row~1 is recovered unchanged, which
confirms that the improvement in row~3 is not an artefact of perturbing the
system.

\subsection{Remaining gap}

Row~3 improves on its baseline but does not attain the machine precision of
row~1. Two effects account for the difference, and both concern the
\emph{solve} rather than the extraction, which \eqref{eq:cross} has repaired
exactly.

First, $\mathcal{X}_i$ is evaluated at $u^*$ and held fixed, whereas
eq.~\eqref{eq:local_solve} solves for $u_i \neq u^*|_{\Omega_i}$; correction and
unknown are therefore mutually inconsistent. Second, as noted above, the
assembly required by \eqref{eq:cross} leaves $A_i$ one-sided but
\emph{under-coupled}: the missing arc weight is compensated on the right-hand
side rather than restored to the operator being inverted.

Removing both would require the cross coupling to enter $A_i$ itself rather than
its right-hand side, which entails either an outer iteration in which
$\mathcal{X}_i$ is refreshed from the neighbour's current iterate, or subdomains
overlapping by the vessel radius so that each local operator carries the full
coupling stencil. We record this as open.

\subsection{Summary}

Eqs.~\eqref{eq:discrete_residual}--\eqref{eq:local_solve} are exact for interfaces
crossed by no vessel, and the failure observed otherwise is neither an
implementation defect nor a discretisation artefact: it is a missing term. In a
3D--1D coupled problem two subdomains interact through two distinct channels ---
the flux across the artificial surface, and the vessel network itself. The
transmission condition as written accounts for the first only. Supplying
\eqref{eq:cross} restores the structural property on which the extraction
depends; embedding that coupling in the local operator remains to be done.
